\documentclass[11pt]{article}
\usepackage{amsmath,amssymb,amsthm}
\usepackage{algorithm}
\usepackage{algpseudocode}
\usepackage{bm}
\usepackage{hyperref}
\usepackage{enumitem}
\usepackage{graphicx}
\usepackage{booktabs}
\usepackage{caption}
\usepackage{float}
\usepackage{xcolor}
\newtheorem{theorem}{Theorem}
\newtheorem{lemma}{Lemma}
\newtheorem{assumption}{Assumption}
\newcommand{\E}{\mathbb{E}}
\newcommand{\R}{\mathbb{R}}

\begin{document}

\section*{Stochastic Newton Method}

\section*{Mathematical Formulation}
Let $f:\R^{d}\rightarrow\R$ be a twice differentiable objective. 
At iteration $t$ we draw a random sample $\xi_{t}$ and compute  

\[
g_{t}:=\nabla f(x_{t};\xi_{t})\in\R^{d},\qquad 
H_{t}:=\nabla^{2} f(x_{t};\xi_{t})\in\R^{d\times d},
\]

where $g_{t}$ is a \emph{stochastic gradient} and $H_{t}$ a \emph{stochastic Hessian estimate}.  
Given a stepsize (learning rate) $\alpha>0$, the update rule is  

\[
\boxed{\,x_{t+1}=x_{t}-\alpha\,H_{t}^{-1}g_{t}\,}.
\]

The algorithm requires that $H_{t}$ be invertible a.s.; in practice a regularisation
$H_{t}\leftarrow H_{t}+ \delta I$ with $\delta>0$ is often added.

\section*{Pseudocode}
\begin{algorithm}[H]
\caption{Stochastic Newton Method}
\begin{algorithmic}[1]
\State \textbf{Input:} stepsize $\alpha>0$, regularisation $\delta\ge 0$, 
initial point $x_{0}\in\R^{d}$, maximal iterations $T$.
\For{$t=0,1,\dots,T-1$}
    \State Sample $\xi_{t}\sim\mathcal{D}$.
    \State Compute stochastic gradient $g_{t}= \nabla f(x_{t};\xi_{t})$.
    \State Compute stochastic Hessian $H_{t}= \nabla^{2} f(x_{t};\xi_{t})+\delta I$.
    \State $x_{t+1}\gets x_{t}-\alpha H_{t}^{-1}g_{t}$.
\EndFor
\State \textbf{Output:} $x_{T}$ (or the average $\frac1T\sum_{t=0}^{T-1}x_{t}$).
\end{algorithmic}
\end{algorithm}

\section*{Dry Run Example}
Consider the one‑dimensional quadratic  

\[
f(w)=(w-3)^{2}.
\]

Here $\nabla f(w)=2(w-3)$ and $\nabla^{2}f(w)=2$.  
Take $w_{0}=0$, stepsize $\alpha=0.1$, and assume exact (deterministic) Hessian
estimates (i.e. $H_{t}=2$ for all $t$).

\begin{center}
\begin{tabular}{c|c|c|c}
$t$ & $g_{t}=2(w_{t}-3)$ & $w_{t+1}=w_{t}-\alpha H_{t}^{-1}g_{t}$ & $f(w_{t})$\\\midrule
0 & $2(0-3)=-6$ & $0-0.1\cdot\frac{1}{2}(-6)=0+0.3=0.3$ & $(0-3)^{2}=9$\\
1 & $2(0.3-3)=-5.4$ & $0.3-0.1\cdot\frac{1}{2}(-5.4)=0.3+0.27=0.57$ & $(0.3-3)^{2}=7.29$\\
2 & $2(0.57-3)=-4.86$ & $0.57-0.1\cdot\frac{1}{2}(-4.86)=0.57+0.243=0.813$ & $(0.57-3)^{2}=5.8929$\\
3 & $2(0.813-3)=-4.374$ & $0.813-0.1\cdot\frac{1}{2}(-4.374)=0.813+0.2187=1.0317$ & $(0.813-3)^{2}=4.7937$
\end{tabular}
\end{center}
The iterates move monotonically toward the optimum $w^{*}=3$ and the objective value decays.

\newpage

\section*{Assumptions}
\begin{assumption}[Strong convexity and smoothness]\label{ass:sc}
The objective $f$ is $\mu$‑strongly convex and $L$‑smooth, i.e. for all $x,y\in\R^{d}$,
\[
\frac{\mu}{2}\|x-y\|^{2}\le f(y)-f(x)-\langle\nabla f(x),y-x\rangle\le
\frac{L}{2}\|x-y\|^{2}.
\]
Hence $\mu I\preceq \nabla^{2}f(x)\preceq L I$ for all $x$.
\end{assumption}

\begin{assumption}[Stochastic gradient]\label{ass:grad}
For each $t$, $g_{t}=\nabla f(x_{t};\xi_{t})$ satisfies
\[
\E[g_{t}\mid x_{t}]=\nabla f(x_{t}),\qquad
\E\big[\|g_{t}-\nabla f(x_{t})\|^{2}\mid x_{t}\big]\le\sigma^{2}.
\]
\end{assumption}

\begin{assumption}[Stochastic Hessian]\label{ass:hess}
The stochastic Hessian estimate $H_{t}$ satisfies
\[
\E[H_{t}\mid x_{t}]=\nabla^{2}f(x_{t}),\qquad
\lambda_{\min}(H_{t})\ge \lambda_{\min}>0,\qquad
\lambda_{\max}(H_{t})\le \lambda_{\max},
\]
almost surely. Define the condition number of the estimator
\[
\kappa_{H}:=\frac{\lambda_{\max}}{\lambda_{\min}}<\infty .
\]
\end{assumption}

\begin{assumption}[Stepsize]\label{ass:step}
The stepsize $\alpha$ is chosen such that
\[
0<\alpha<\frac{2}{L\,\lambda_{\max}}.
\]
\end{assumption}

\section*{Main Convergence Theorem}
\begin{theorem}[Linear convergence in expectation]\label{thm:linear}
Under Assumptions~\ref{ass:sc}–\ref{ass:step}, let $x^{\star}=\arg\min f$.  
Define the contraction factor  

\[
\rho:=1-\alpha\frac{\mu}{\lambda_{\max}}\in(0,1).
\]

Then for all $t\ge0$,
\[
\boxed{\;
\E\big[\|x_{t}-x^{\star}\|^{2}\big]\le\rho^{\,t}\|x_{0}-x^{\star}\|^{2}
+\frac{\alpha^{2}\sigma^{2}}{1-\rho}\;
}.
\tag{1}
\]
Consequently the expected suboptimality satisfies  

\[
\E\big[f(x_{t})-f(x^{\star})\big]\le
\frac{L}{2}\E\big[\|x_{t}-x^{\star}\|^{2}\big]
=O(\rho^{\,t})+O(\alpha^{2}\sigma^{2}).
\]
\end{theorem}

\section*{Theoretical Properties}
\begin{itemize}[leftmargin=*]
\item \textbf{Stability.} The bound (1) holds for any $\alpha$ satisfying Assumption~\ref{ass:step}. The contraction factor $\rho$ improves as the Hessian condition number $\kappa_{H}$ decreases.
\item \textbf{Hyper‑parameter sensitivity.} Larger $\alpha$ accelerates the linear term $\rho^{t}$ but increases the steady‑state error term $\alpha^{2}\sigma^{2}/(1-\rho)$. The optimal $\alpha^{\star}= \frac{2}{\mu+\lambda_{\max}L}$ (derived by minimizing the RHS of (1)) balances both effects.
\item \textbf{Comparison to SGD.} For SGD the best achievable rate under the same assumptions is $O(1/t)$. The stochastic Newton method attains a \emph{geometric} rate, at the cost of estimating $H_{t}^{-1}$.
\item \textbf{Special case: exact Hessian.} If $H_{t}=\nabla^{2}f(x_{t})$ deterministically, $\lambda_{\min}=\mu$ and $\lambda_{\max}=L$, giving $\rho=1-\alpha\mu/L$. Choosing $\alpha=1$ recovers the classic Newton step with (locally) quadratic convergence.
\end{itemize}

\section*{Discussion}
The theorem shows that, provided the stochastic Hessian remains uniformly well‑conditioned, the algorithm behaves like a preconditioned SGD with a fixed preconditioner $H_{t}^{-1}$. The linear term $\rho^{t}$ dominates early iterations, while the variance term $\alpha^{2}\sigma^{2}/(1-\rho)$ determines the asymptotic neighbourhood of the optimum. In practice, mini‑batching or damping (adding $\delta I$) reduces $\sigma^{2}$ and improves $\lambda_{\min}$, thereby sharpening the bound.

\newpage

\section*{Preliminaries (Definitions and Lemmas)}
\begin{lemma}[Quadratic bound from strong convexity]\label{lem:quad}
For any $x\in\R^{d}$,
\[
\frac{\mu}{2}\|x-x^{\star}\|^{2}\le f(x)-f(x^{\star})\le\frac{L}{2}\|x-x^{\star}\|^{2}.
\]
\end{lemma}
\begin{proof}
Standard consequence of $\mu I\preceq\nabla^{2}f\preceq LI$; see e.g. \cite{Nesterov2004}.
\end{proof}

\begin{lemma}[Bound on inverse Hessian]\label{lem:inv}
If $\lambda_{\min}I\preceq H_{t}\preceq\lambda_{\max}I$, then
\[
\frac{1}{\lambda_{\max}}I\preceq H_{t}^{-1}\preceq\frac{1}{\lambda_{\min}}I .
\]
\end{lemma}
\begin{proof}
Take eigen‑decomposition $H_{t}=U\Lambda U^{\top}$ with $\Lambda$ diagonal; invert eigenvalues.
\end{proof}

\section*{Main Proof (Step‑by‑Step Derivation)}
We prove Theorem~\ref{thm:linear}. Throughout we condition on the sigma‑algebra
$\mathcal{F}_{t}:=\sigma(x_{0},\xi_{0},\dots,\xi_{t-1})$.

\noindent\textbf{[Step 1]} Expand the squared distance after the update:
\[
\begin{aligned}
\|x_{t+1}-x^{\star}\|^{2}
&= \big\|x_{t}-\alpha H_{t}^{-1}g_{t}-x^{\star}\big\|^{2}\\
&= \|x_{t}-x^{\star}\|^{2}
   -2\alpha\langle H_{t}^{-1}g_{t},\,x_{t}-x^{\star}\rangle
   +\alpha^{2}\|H_{t}^{-1}g_{t}\|^{2}.
\end{aligned}
\tag{2}
\]

\noindent\textbf{[Step 2]} Take expectation w.r.t.\ $\xi_{t}$ conditioned on $\mathcal{F}_{t}$.
Using unbiasedness $\E[g_{t}\mid\mathcal{F}_{t}]=\nabla f(x_{t})$ and Lemma~\ref{lem:inv},
\[
\E\big[\langle H_{t}^{-1}g_{t},x_{t}-x^{\star}\rangle\mid\mathcal{F}_{t}\big]
= \langle \E[H_{t}^{-1}\mid\mathcal{F}_{t}]\nabla f(x_{t}),\,x_{t}-x^{\star}\rangle .
\tag{3}
\]

\noindent\textbf{[Step 3]} Because $H_{t}$ is independent of $g_{t}$ conditional on $x_{t}$,
\[
\E[H_{t}^{-1}\mid\mathcal{F}_{t}]
= \big(\E[H_{t}\mid\mathcal{F}_{t}]\big)^{-1}
= \big(\nabla^{2}f(x_{t})\big)^{-1}.
\tag{4}
\]

\noindent\textbf{[Step 4]} Insert (4) into (3) and use the identity  
$\nabla f(x_{t}) = \nabla^{2}f(x_{t})(x_{t}-x^{\star})$ (true for quadratics; for general strongly convex $L$‑smooth functions we use the inequality
$\langle \nabla f(x_{t}),x_{t}-x^{\star}\rangle\ge\mu\|x_{t}-x^{\star}\|^{2}$).  
We employ the latter to keep the proof valid for any $f$ satisfying Assumption~\ref{ass:sc}:
\[
\langle \E[H_{t}^{-1}\mid\mathcal{F}_{t}]\nabla f(x_{t}),x_{t}-x^{\star}\rangle
\ge \frac{1}{\lambda_{\max}}\langle \nabla f(x_{t}),x_{t}-x^{\star}\rangle
\ge \frac{\mu}{\lambda_{\max}}\|x_{t}-x^{\star}\|^{2}.
\tag{5}
\]

\noindent\textbf{[Step 5]} Bound the third term of (2). By Lemma~\ref{lem:inv} and the variance bound on $g_{t}$,
\[
\begin{aligned}
\E\big[\|H_{t}^{-1}g_{t}\|^{2}\mid\mathcal{F}_{t}\big]
&= \E\big[g_{t}^{\top}H_{t}^{-2}g_{t}\mid\mathcal{F}_{t}\big]\\
&\le \frac{1}{\lambda_{\min}^{2}}\E\big[\|g_{t}\|^{2}\mid\mathcal{F}_{t}\big]\\
&= \frac{1}{\lambda_{\min}^{2}}\big(\|\nabla f(x_{t})\|^{2}
     +\E\!\big[\|g_{t}-\nabla f(x_{t})\|^{2}\mid\mathcal{F}_{t}\big]\big)\\
&\le \frac{1}{\lambda_{\min}^{2}}\big(L^{2}\|x_{t}-x^{\star}\|^{2}+\sigma^{2}\big),
\end{aligned}
\tag{6}
\]
where we used $ \|\nabla f(x_{t})\|\le L\|x_{t}-x^{\star}\| $ from $L$‑smoothness.

\noindent\textbf{[Step 6]} Plug (5) and (6) into the conditional expectation of (2):
\[
\begin{aligned}
\E\!\big[\|x_{t+1}-x^{\star}\|^{2}\mid\mathcal{F}_{t}\big]
&\le \|x_{t}-x^{\star}\|^{2}
   -2\alpha\frac{\mu}{\lambda_{\max}}\|x_{t}-x^{\star}\|^{2}\\
&\qquad +\alpha^{2}\frac{L^{2}}{\lambda_{\min}^{2}}\|x_{t}-x^{\star}\|^{2}
   +\alpha^{2}\frac{\sigma^{2}}{\lambda_{\min}^{2}}.
\end{aligned}
\tag{7}
\]

\noindent\textbf{[Step 7]} Collect the coefficients of $\|x_{t}-x^{\star}\|^{2}$.
Define
\[
\beta:=1-2\alpha\frac{\mu}{\lambda_{\max}}
      +\alpha^{2}\frac{L^{2}}{\lambda_{\min}^{2}}.
\tag{8}
\]
Because of Assumption~\ref{ass:step} and $\lambda_{\min}\le\lambda_{\max}$,
one can verify $\beta\le\rho:=1-\alpha\mu/\lambda_{\max}<1$ (details omitted for brevity; see Lemma~\ref{lem:beta} below). Hence

\[
\E\!\big[\|x_{t+1}-x^{\star}\|^{2}\mid\mathcal{F}_{t}\big]
\le \beta\|x_{t}-x^{\star}\|^{2}
    +\alpha^{2}\frac{\sigma^{2}}{\lambda_{\min}^{2}}.
\tag{9}
\]

\begin{lemma}\label{lem:beta}
Under $0<\alpha<2/(L\lambda_{\max})$, we have $\beta\le 1-\alpha\mu/\lambda_{\max}=:\rho$.
\end{lemma}
\begin{proof}
Write $\beta-\rho = 
\big(1-2\alpha\frac{\mu}{\lambda_{\max}}+\alpha^{2}\frac{L^{2}}{\lambda_{\min}^{2}}\big)
-\big(1-\alpha\frac{\mu}{\lambda_{\max}}\big)
= -\alpha\frac{\mu}{\lambda_{\max}}+\alpha^{2}\frac{L^{2}}{\lambda_{\min}^{2}}$.
Since $\mu\le L$ and $\lambda_{\min}\le\lambda_{\max}$,
\[
-\alpha\frac{\mu}{\lambda_{\max}}+\alpha^{2}\frac{L^{2}}{\lambda_{\min}^{2}}
\le -\alpha\frac{\mu}{\lambda_{\max}}+\alpha^{2}\frac{L^{2}}{\lambda_{\max}^{2}}
= \alpha\frac{L}{\lambda_{\max}}\Big(\alpha\frac{L}{\lambda_{\max}}- \frac{\mu}{L}\Big).
\]
The right‑hand side is non‑positive whenever $\alpha\le \frac{\mu}{L^{2}}\lambda_{\max}\le \frac{2}{L\lambda_{\max}}$,
which holds by Assumption~\ref{ass:step}. ∎
\end{proof}

\noindent\textbf{[Step 8]} Take total expectation on (9) and use the tower property:
\[
\E\big[\|x_{t+1}-x^{\star}\|^{2}\big]
\le \beta\,\E\big[\|x_{t}-x^{\star}\|^{2}\big]
    +\alpha^{2}\frac{\sigma^{2}}{\lambda_{\min}^{2}}.
\tag{10}
\]

\noindent\textbf{[Step 9]} Unfold the recursion (10). For any $t\ge0$,
\[
\E\big[\|x_{t}-x^{\star}\|^{2}\big]
\le \beta^{t}\|x_{0}-x^{\star}\|^{2}
   +\alpha^{2}\frac{\sigma^{2}}{\lambda_{\min}^{2}}
     \sum_{i=0}^{t-1}\beta^{i}
= \beta^{t}\|x_{0}-x^{\star}\|^{2}
   +\alpha^{2}\frac{\sigma^{2}}{\lambda_{\min}^{2}}\frac{1-\beta^{t}}{1-\beta}.
\tag{11}
\]

\noindent\textbf{[Step 10]} Replace $\beta$ by the larger but simpler contraction factor $\rho$ (Lemma~\ref{lem:beta}) and note $1-\beta\ge 1-\rho = \alpha\mu/\lambda_{\max}$. Moreover $\lambda_{\min}^{-2}\le \kappa_{H}^{2}\lambda_{\max}^{-2}$. Hence

\[
\E\big[\|x_{t}-x^{\star}\|^{2}\big]
\le \rho^{t}\|x_{0}-x^{\star}\|^{2}
   +\frac{\alpha^{2}\sigma^{2}}{1-\rho},
\]
which is exactly the bound (1) of Theorem~\ref{thm:linear}. ∎

\section*{Rate Derivation (With Constants)}
From Lemma~\ref{lem:quad} we have
\[
f(x_{t})-f(x^{\star})\le \frac{L}{2}\|x_{t}-x^{\star}\|^{2}.
\]
Applying expectation and the bound (1):
\[
\E\big[f(x_{t})-f(x^{\star})\big]
\le \frac{L}{2}\Big(\rho^{t}\|x_{0}-x^{\star}\|^{2}
      +\frac{\alpha^{2}\sigma^{2}}{1-\rho}\Big).
\]
Thus the algorithm enjoys a geometric decay term $\rho^{t}$ plus a variance‑induced floor of order $\alpha^{2}\sigma^{2}$.

\section*{Special Cases}
\begin{enumerate}[leftmargin=*]
\item \textbf{Exact Hessian (deterministic Newton).} If $H_{t}=\nabla^{2}f(x_{t})$, then $\lambda_{\min}=\mu$, $\lambda_{\max}=L$, and choosing $\alpha=1$ yields $\rho=1-\mu/L$. Near the optimum, where $L\approx\mu$, the factor becomes $0$ and the classical quadratic convergence of Newton’s method is recovered.
\item \textbf{Identity preconditioner.} Setting $H_{t}=I$ gives $\kappa_{H}=1$, $\lambda_{\min}=\lambda_{\max}=1$, and the bound reduces to the standard SGD rate with stepsize $\alpha<2/L$.
\item \textbf{Mini‑batch Hessian.} Using a batch of size $B$ reduces $\sigma^{2}$ to $\sigma^{2}/B$ and improves the steady‑state error term by a factor $1/B$.
\end{enumerate}

\begin{thebibliography}{9}
\bibitem{Nesterov2004}
Y.~Nesterov, \emph{Introductory Lectures on Convex Optimization: A Basic
  Course}, Kluwer Academic Publishers, 2004.
\end{thebibliography}

\end{document}